\chapter{Quadratic forms and quadrics}\label{geometry:ch:quadrics}

\begin{center}\itshape
A demonstration of the theorem that every homogeneous quadratic polynomial is reducible by real orthogonal substitutions\dots
\upshape(J.~J.~Sylvester, title of a paper in the \emph{Philosophical Magazine}, 1852)
\end{center}

In 1896 the All-Russian Industrial and Art Exhibition opened in Nizhny Novgorod, and its visitors could see a water tower unlike any other. It
was curved like an hourglass and yet it was made entirely of straight steel bars, crossing each other in two directions like the strings of a
lampshade. Its engineer, Vladimir Shukhov, had patented the idea. The surface is a \emph{hyperboloid of one sheet}, and it has a remarkable property:
through every one of its points pass two straight lines lying entirely on the surface. Straight bars are cheap, light and strong, and Shukhov built
dozens of such towers. The most famous, the radio tower on Shabolovka street in Moscow, six hyperboloids stacked on top of each other and 150 metres
high, was finished in March 1922.

The mathematics behind the tower is the algebra of quadratic forms. A quadric surface is the zero set of a polynomial of degree two in three
variables, and its shape is decided by a few numbers attached to a symmetric matrix. James Joseph Sylvester found the most important of them in
1852: however you change coordinates, a real quadratic form always has the same number of positive and the same number of negative squares. He
called this the law of inertia. This chapter proves it, uses it to classify the quadric surfaces, explains the two families of lines on the
hyperboloid, and then leaves the real numbers for arbitrary fields, where Witt's cancellation theorem replaces Sylvester's law, and for finite fields,
where quadrics can be counted point by point.

\begin{think}
The surface $z=xy$ is a saddle. Is it curved? Try to find a straight line lying on it. How many lines pass through the origin? Through the point
$(1,1,1)$?
\end{think}

\section{Symmetric bilinear forms}

Let $K$ be a field with $1+1\ne0$ and $V$ a vector space of dimension $n$ over $K$. A \emph{quadratic form} on $V$ is a function $q(v)=\beta(v,v)$,
where $\beta$ is a symmetric bilinear form; since $2\beta(u,v)=q(u+v)-q(u)-q(v)$, the form $q$ determines $\beta$. In a basis $e_1,\dots,e_n$ the
form is described by its \emph{Gram matrix} $A=(\beta(e_i,e_j))$, so that $q(x)=x^{\mathsf T}Ax$ in coordinates. If $P$ is the matrix of a change of
basis, the new Gram matrix is $P^{\mathsf T}AP$. Two forms are \emph{equivalent}, or \emph{isometric}, if their matrices are related in this way,
that is, if some linear isomorphism carries one to the other.

The rank of $A$ does not change under $A\mapsto P^{\mathsf T}AP$, so it is called the rank of $q$. The \emph{radical} $V^\perp=\{v:\beta(v,w)=0\text{ for all }w\}$
is the kernel of $A$, and $q$ is \emph{nondegenerate} when the radical is $0$, that is, when $\det A\ne0$. For a subspace $W$ we write
$W^\perp=\{v:\beta(v,w)=0\text{ for all }w\in W\}$, and $U\perp W$ means that $\beta(U,W)=0$. If $q$ is nondegenerate on $W$ (the restriction has
nonzero determinant), then $V=W\oplus W^\perp$: the map $v\mapsto\beta(v,\cdot)|_W$ from $V$ to the dual of $W$ is onto (already its restriction to
$W$ is an isomorphism), so $\dim W^\perp=n-\dim W$, and $W\cap W^\perp=0$ by nondegeneracy on $W$.

\begin{proposition}[Orthogonal bases]\label{geometry:prop:orthbasis}
Every quadratic form has an orthogonal basis, that is, a basis in which its matrix is diagonal: $q=a_1x_1^2+\dots+a_nx_n^2$.
\end{proposition}

\begin{proof}
By induction on $n$. If $q=0$ then $\beta=0$ by the polarisation formula, and every basis is orthogonal. Otherwise choose $e_1$ with $q(e_1)\ne0$.
The form is nondegenerate on the line $Ke_1$, so $V=Ke_1\oplus e_1^\perp$. An orthogonal basis of $e_1^\perp$, which exists by induction, together
with $e_1$ is an orthogonal basis of $V$.
\end{proof}

We write $\langle a_1,\dots,a_n\rangle$ for the diagonal form $a_1x_1^2+\dots+a_nx_n^2$. The coefficients are not unique: replacing $e_i$ by $\lambda
e_i$ multiplies $a_i$ by $\lambda^2$. Over $\mathbb C$ every nonzero $a_i$ can be made $1$, so complex forms are classified by their rank alone. Over
$\mathbb R$ every nonzero $a_i$ can be made $\pm1$, and the real question is how many of each sign there are.

\section{Sylvester's law of inertia}

\begin{theorem}[Sylvester's law of inertia]\label{geometry:thm:sylvester}
Let $q$ be a real quadratic form with an orthogonal basis in which $q=x_1^2+\dots+x_p^2-x_{p+1}^2-\dots-x_{p+m}^2$. Then $p$ is the largest dimension
of a subspace on which $q$ is positive definite, and $m$ the largest dimension of a subspace on which $q$ is negative definite. In particular $p$ and
$m$ do not depend on the basis.
\end{theorem}

\begin{proof}
On $W_0=\langle e_1,\dots,e_p\rangle$, $q(x)=x_1^2+\dots+x_p^2>0$ for $x\ne0$, so the largest dimension is at least $p$. Conversely, let $W$ be a
subspace on which $q$ is positive definite, and let $N=\langle e_{p+1},\dots,e_n\rangle$. On $N$ we have $q=-x_{p+1}^2-\dots-x_{p+m}^2\le0$, so
$W\cap N=0$, and $\dim W\le n-\dim N=p$. The statement for $m$ follows by applying this to $-q$.
\end{proof}

The pair $(p,m)$ is the \emph{signature} of $q$, and $p+m$ is its rank. So two real forms are equivalent if and only if they have the same signature.
The spectral theorem of Volume~2 gives a second way to read the signature: in an orthonormal basis of eigenvectors of $A$, the form is
$\sum\lambda_ix_i^2$, so $p$ and $m$ are the numbers of positive and negative eigenvalues of $A$, counted with multiplicity.

\begin{figure}[ht]
\centering
\fig{geometry-ch06-lorentz}
\caption{The level sets of the form $x^2+y^2-t^2$, of signature $(2,1)$: the value $1$ gives a hyperboloid of one sheet, $0$ the light cone and
$-1$ a hyperboloid of two sheets.}
\label{geometry:fig:lorentz}
\end{figure}

\begin{corollary}\label{geometry:cor:realconics}
A nondegenerate real projective conic is either empty or projectively equivalent to the circle $x^2+y^2=z^2$. A nondegenerate real projective quadric
surface in $\mathbb P^3(\mathbb R)$ is empty, or equivalent to the sphere $x^2+y^2+z^2=w^2$, or to $x^2+y^2=z^2+w^2$.
\end{corollary}

\begin{proof}
Multiplying the form by $-1$ does not change the zero set, so we may assume $p\ge m$. For conics the nondegenerate signatures are then $(3,0)$, with
no real zero, and $(2,1)$. For surfaces they are $(4,0)$, $(3,1)$ and $(2,2)$.
\end{proof}

The two nonempty surfaces are really different: the sphere contains no line, because a line would be a two-dimensional subspace on which the
form of signature $(3,1)$ vanishes, and it would meet the three-dimensional subspace where the form is positive definite. The quadric of
signature $(2,2)$ is full of lines, as we will see in Section~\ref{geometry:sec:ruled}.

\section{Quadric surfaces in space}

A \emph{quadric surface} in $\mathbb R^3$ is the zero set of a polynomial of degree two,
\[
x^{\mathsf T}Mx+2b^{\mathsf T}x+c=0,\qquad
\hat A=\begin{pmatrix}M&b\\b^{\mathsf T}&c\end{pmatrix},
\]
with $M$ a nonzero symmetric $3\times3$ matrix. The $4\times4$ matrix $\hat A$ is the matrix of the homogenised form in $\mathbb P^3(\mathbb R)$, and
the quadric is \emph{nondegenerate} when $\det\hat A\ne0$.

\begin{theorem}\label{geometry:thm:quadrics}
A nonempty nondegenerate quadric surface in $\mathbb R^3$ can be brought by an affine change of coordinates to exactly one of the forms
\[
\begin{gathered}
x^2+y^2+z^2=1,\qquad x^2+y^2-z^2=1,\qquad x^2-y^2-z^2=1,\\
z=x^2+y^2,\qquad z=x^2-y^2,
\end{gathered}
\]
the ellipsoid, the hyperboloids of one and of two sheets, and the elliptic and hyperbolic paraboloids.
\end{theorem}

\begin{proof}
\emph{Case 1: $M$ invertible.} The translation $x\mapsto x-M^{-1}b$ removes the linear term and turns the equation into $x^{\mathsf T}Mx=c'$ with
$c'=b^{\mathsf T}M^{-1}b-c$. Subtracting from the last row of $\hat A$ the combination $b^{\mathsf T}M^{-1}$ of the first three rows shows
$\det\hat A=\det M\cdot(c-b^{\mathsf T}M^{-1}b)$, so $c'\ne0$. Dividing by $c'$ and diagonalising the form $x^{\mathsf T}Mx/c'$ with Sylvester's law gives
$\pm x^2\pm y^2\pm z^2=1$. Up to the order of the variables there are four sign patterns: three plus signs (ellipsoid), two (one sheet), one (two
sheets), none (empty).

\emph{Case 2: $M$ singular.} The first three columns of $\hat A$ are linearly independent, and they are the columns of $M$ with one more entry each,
so $\operatorname{rank}M\ge2$. Hence $\operatorname{rank}M=2$, and after a linear change of coordinates $M=\mathrm{diag}(d_1,d_2,0)$ with $d_1,d_2\ne0$.
The equation reads $d_1x^2+d_2y^2+2b_1x+2b_2y+2b_3z+c=0$. Completing the squares in $x$ and $y$ removes $b_1$ and $b_2$. Then $b_3\ne0$, since otherwise
the third row of $\hat A$ would be zero. The substitution $z\mapsto-(z+c'')/(2b_3)$, for the appropriate constant $c''$, gives $d_1x^2+d_2y^2=z$, and
rescaling $x$ and $y$ gives $z=\pm x^2\pm y^2$. Changing $z$ to $-z$ if needed, this is $z=x^2+y^2$ or $z=x^2-y^2$.

The five surfaces are not affinely equivalent. An affine change of coordinates acts on $M$ by $M\mapsto P^{\mathsf T}MP$, and we may also multiply
the equation by a nonzero constant, so the rank of $M$ and its signature up to order are invariants. They are $(3,0)$ for the ellipsoid, $(2,1)$ for
both hyperboloids, and $(2,0)$ and $(1,1)$ for the two paraboloids. The hyperboloids are told apart by their connected pieces: the one-sheeted
hyperboloid is connected, while $x^2-y^2-z^2=1$ has the two pieces $x\ge1$ and $x\le-1$.
\end{proof}

\begin{remark}
Counting also the degenerate cases (cones, cylinders, pairs of planes, lines, points) and the ``imaginary'' ones with no real points, the classical
list of normal forms of real quadrics has seventeen entries; see Berger's \emph{Geometry II}, Chapter~15.
\end{remark}

In projective terms, the ellipsoid, the two-sheeted hyperboloid and the elliptic paraboloid are the same surface, the projective sphere of
signature $(3,1)$, meeting the plane at infinity in nothing, in a conic, or in a single point where it is tangent. The hyperboloid of one sheet and
the hyperbolic paraboloid are the ruled quadric of signature $(2,2)$, cut by the plane at infinity in a conic or in two lines.

\section{Ruled quadrics}\label{geometry:sec:ruled}

The form $x^2+y^2-z^2-w^2$ can be rewritten: with $x_0=x-z$, $x_3=x+z$, $x_1=w-y$, $x_2=w+y$ we get $x_0x_3-x_1x_2=x^2-z^2-w^2+y^2$. So the quadric of
signature $(2,2)$ is projectively the quadric
\[
Q=\{[x_0:x_1:x_2:x_3]\in\mathbb P^3:\ x_0x_3=x_1x_2\},
\]
and we can study it over any field. Write a point of $\mathbb P^3$ as a nonzero $2\times2$ matrix $\begin{pmatrix}x_0&x_1\\x_2&x_3\end{pmatrix}$ up to a
factor. Then $Q$ is the set of matrices of determinant $0$, that is, of rank $1$.

\begin{theorem}[The Segre quadric]\label{geometry:thm:segre}
The map $\sigma:\mathbb P^1\times\mathbb P^1\to Q$, $\sigma([s:t],[u:v])=[su:sv:tu:tv]$, is a bijection. For each fixed $a\in\mathbb P^1$, the sets
$\sigma(\{a\}\times\mathbb P^1)$ and $\sigma(\mathbb P^1\times\{a\})$ are lines. Two lines of the same family are disjoint, two lines of different families
meet in exactly one point, and every line contained in $Q$ belongs to one of the two families.
\end{theorem}

\begin{proof}
The point $\sigma([s:t],[u:v])$ is the matrix $\binom st(u\ v)$, a column times a row. A matrix of rank $1$ has a one-dimensional column space spanned
by some $\binom st$, so it is $\binom st(u\ v)$ for some nonzero row; the column $\binom st$ and the row $(u\ v)$ are determined up to factors by
the matrix (its column space and its row space). So $\sigma$ is a bijection. For fixed $[s:t]$, the matrices $\binom st(u\ v)$ form a two-dimensional
subspace, a line of $\mathbb P^3$; likewise for fixed $[u:v]$. Two lines of the same family come from different columns (or rows), so by bijectivity
they are disjoint; the lines $\{[s:t]\}\times\mathbb P^1$ and $\mathbb P^1\times\{[u:v]\}$ meet exactly in $\sigma([s:t],[u:v])$.

Let $L\subset Q$ be a line, a two-dimensional space of matrices of rank at most one. Take two independent elements $c_1r_1$ and $c_2r_2$ (columns
times rows). If $c_1,c_2$ were independent and $r_1,r_2$ were independent, then $c_1r_1+c_2r_2$ would map the independent vectors dual to $r_1,r_2$
onto independent vectors, so it would have rank $2$. So the columns are proportional, and then $L$ lies in a line of the first family (and equals it,
both being lines), or the rows are proportional and $L$ is a line of the second family.
\end{proof}

\begin{figure}[ht]
\centering
\fig{geometry-ch06-hyperboloid}
\caption{The hyperboloid of one sheet $x^2+y^2-z^2=1$ with its two families of lines. Through $(\cos\theta,\sin\theta,0)$ pass the lines with directions
$(-\sin\theta,\cos\theta,\pm1)$.}
\label{geometry:fig:hyperboloid}
\end{figure}

On the hyperboloid $x^2+y^2-z^2=1$ the lines are explicit: through each point $(\cos\theta,\sin\theta,0)$ of the waist pass the lines
$(\cos\theta-s\sin\theta,\ \sin\theta+s\cos\theta,\ \pm s)$, and indeed $(\cos\theta-s\sin\theta)^2+(\sin\theta+s\cos\theta)^2=1+s^2$. These are
Shukhov's bars. On the saddle $z=xy$ they are even simpler, the lines $x=c$ and $y=c$ (Figure~\ref{geometry:fig:saddle}).

\begin{figure}[ht]
\centering
\fig{geometry-ch06-saddle}
\caption{The saddle $z=xy$ is ruled twice: by the lines $x=c,\ z=cy$ (charcoal) and $y=c,\ z=cx$ (rust).}
\label{geometry:fig:saddle}
\end{figure}

\begin{beauty}
Three pairwise skew lines in space determine a unique quadric surface, and it is ruled: Exercise~\ref{geometry:ex:skewlines}. Every line that meets
the three lines lies on it, and these lines form the second family. So the hourglass shape of Shukhov's tower is forced by any three of its bars.
\end{beauty}

\section{Isotropic vectors and Witt's theorem}

Over a general field, the sign of a number no longer makes sense, and Sylvester's law is replaced by a cancellation theorem. A nonzero vector $v$ is
\emph{isotropic} if $q(v)=0$, and $q$ is \emph{anisotropic} if it has no isotropic vector. The \emph{hyperbolic plane} $\mathbb H$ is the form $2xy$
on $K^2$, with Gram matrix $\begin{pmatrix}0&1\\1&0\end{pmatrix}$. It is also $\langle1,-1\rangle$: with $x=X-Y$ and $y=(X+Y)/2$ we get
$2xy=X^2-Y^2$.

\begin{lemma}\label{geometry:lem:hyperbolic}
Let $q$ be nondegenerate and $e$ an isotropic vector. Then $e$ lies in a hyperbolic plane $H$, and $V=H\oplus H^\perp$ with $H^\perp$ nondegenerate.
\end{lemma}

\begin{proof}
Since $e$ is not in the radical, there is $f'$ with $\beta(e,f')=1$. Let $f=f'-\frac{q(f')}2e$. Then $\beta(e,f)=1$ and
$q(f)=q(f')-q(f')\beta(e,f')=0$. In the basis $e,f$ the Gram matrix is $\begin{pmatrix}0&1\\1&0\end{pmatrix}$, so $H=\langle e,f\rangle$ is a hyperbolic plane,
nondegenerate, and $V=H\oplus H^\perp$. A vector of the radical of $H^\perp$ would be orthogonal to $H$ and to $H^\perp$, hence to $V$, hence $0$.
\end{proof}

The key tool is the reflection. For $w$ with $q(w)\ne0$, the map
\[
\tau_w(x)=x-\frac{2\beta(x,w)}{q(w)}\,w
\]
fixes $w^\perp$, sends $w$ to $-w$, and preserves $q$: $q(\tau_wx)=q(x)-4\frac{\beta(x,w)^2}{q(w)}+4\frac{\beta(x,w)^2}{q(w)^2}q(w)=q(x)$.

\begin{lemma}\label{geometry:lem:transitive}
If $q(u)=q(v)\ne0$, there is an isometry of $V$, a product of at most two reflections, that maps $u$ to $v$.
\end{lemma}

\begin{proof}
We have $q(u+v)+q(u-v)=2q(u)+2q(v)=4q(u)\ne0$, so $q(u-v)\ne0$ or $q(u+v)\ne0$. If $q(u-v)\ne0$: since $q(u-v)=2q(u)-2\beta(u,v)=2\beta(u,u-v)$, we get
$\tau_{u-v}(u)=u-(u-v)=v$. If $q(u+v)\ne0$: in the same way $\tau_{u+v}(u)=-v$, and then $\tau_v\tau_{u+v}(u)=v$.
\end{proof}

\begin{theorem}[Witt cancellation]\label{geometry:thm:witt}
Let $q$, $q_1$, $q_2$ be quadratic forms with $q$ nondegenerate. If $q\perp q_1\cong q\perp q_2$, then $q_1\cong q_2$.
\end{theorem}

\begin{proof}
Diagonalise $q=\langle a_1,\dots,a_k\rangle$ with all $a_i\ne0$. Cancelling one $\langle a_i\rangle$ at a time, it suffices to treat $q=\langle a\rangle$
with $a\ne0$. Transporting by the given isometry, we have one space $V$ with two decompositions $V=Ku\perp W_1=Kv\perp W_2$, where $q(u)=q(v)=a$ and
$W_i\cong q_i$. Since $a\ne0$, $W_1=u^\perp$ and $W_2=v^\perp$. By Lemma~\ref{geometry:lem:transitive} there is an isometry $\sigma$ of $V$ with $\sigma(u)=v$, and
then $\sigma(W_1)=\sigma(u^\perp)=v^\perp=W_2$. So $q_1\cong q_2$.
\end{proof}

\begin{corollary}[Witt decomposition]\label{geometry:cor:wittindex}
Every nondegenerate form is isometric to $m\mathbb H\perp q_0$ with $q_0$ anisotropic, and the number $m$, the \emph{Witt index}, and the class of $q_0$
are unique.
\end{corollary}

\begin{proof}
Existence: if $q$ is isotropic, split off a hyperbolic plane with Lemma~\ref{geometry:lem:hyperbolic} and continue with $H^\perp$; stop when what remains is
anisotropic. Uniqueness: if $m\mathbb H\perp q_0\cong m'\mathbb H\perp q_0'$ with $m\le m'$, cancelling $m\mathbb H$ gives $q_0\cong(m'-m)\mathbb H\perp q_0'$.
As $q_0$ is anisotropic and $\mathbb H$ is not, $m'=m$, and then $q_0\cong q_0'$.
\end{proof}

Over $\mathbb R$, the Witt index of a nondegenerate form of signature $(p,m)$ is $\min(p,m)$, and the anisotropic part is definite. Over
$\mathbb Q$ the anisotropic parts are a deep subject: whether a rational form is isotropic is decided by the theorem of Hasse and Minkowski, which says
that it is isotropic over $\mathbb Q$ if and only if it is isotropic over $\mathbb R$ and over every $p$-adic field $\mathbb Q_p$. Volume~12 returns to
this question for conics, with Legendre's theorem.

\section{Quadrics over finite fields}

Over a finite field $\mathbb F_q$ with $q$ odd, everything can be counted. The first fact is that small anisotropic forms do not exist.

\begin{proposition}\label{geometry:prop:chevalley}
Every nondegenerate quadratic form in at least three variables over $\mathbb F_q$ ($q$ odd) is isotropic.
\end{proposition}

\begin{proof}
It is enough to treat $ax^2+by^2+cz^2$ with $abc\ne0$. The squares of $\mathbb F_q$ are $0$ and $(q-1)/2$ nonzero elements (the map $x\mapsto x^2$ on
$\mathbb F_q^\times$ is two-to-one). So $\{ax^2:x\in\mathbb F_q\}$ and $\{-c-by^2:y\in\mathbb F_q\}$ both have $(q+1)/2$ elements, and two subsets of
$\mathbb F_q$ of total size $q+1$ must meet: $ax^2=-c-by^2$ for some $x,y$, and $(x,y,1)$ is isotropic.
\end{proof}

So a nondegenerate form in four variables is $\mathbb H\perp q_2$ with $q_2$ a nondegenerate plane, and either $q_2\cong\mathbb H$ (Witt index $2$, the
\emph{hyperbolic} quadric) or $q_2$ is anisotropic (Witt index $1$, the \emph{elliptic} quadric).

\begin{theorem}\label{geometry:thm:countquadric}
A nondegenerate quadric surface in $\mathbb P^3(\mathbb F_q)$ has $(q+1)^2$ points if it is hyperbolic and $q^2+1$ points if it is elliptic.
\end{theorem}

\begin{proof}
Write the form as $xy+N(z,w)$, where $xy$ is a hyperbolic plane (scale one variable to replace $2xy$ by $xy$) and $N$ is the second plane. We count
the vectors $(x,y,z,w)$ with $xy=-N(z,w)$. The equation $xy=0$ has $2q-1$ solutions, and $xy=c$ with $c\ne0$ has $q-1$.

If $N$ is anisotropic, $N(z,w)=0$ only for $(z,w)=0$, so the count is $(2q-1)+(q^2-1)(q-1)$. If $N=zw$, then $zw=0$ for $2q-1$ pairs and $zw\ne0$ for
$(q-1)^2$ pairs, so the count is $(2q-1)^2+(q-1)^2(q-1)$.

Subtract the zero vector and divide by the $q-1$ nonzero scalars. In the elliptic case, $\frac{2q-2+(q^2-1)(q-1)}{q-1}=2+q^2-1=q^2+1$. In the hyperbolic case,
$\frac{(2q-1)^2-1}{q-1}+(q-1)^2=4q+q^2-2q+1=(q+1)^2$.
\end{proof}

The hyperbolic count is no surprise: the quadric is the Segre image of $\mathbb P^1\times\mathbb P^1$, with $(q+1)^2$ points, and it carries
$2(q+1)$ lines. The elliptic quadric carries no line at all (Exercise~\ref{geometry:ex:linesFq}): its $q^2+1$ points are an \emph{ovoid}, a finite analogue of
the sphere, no three of whose points are collinear.

\section{Exercises}

\begin{exercise}[$\star$]\label{geometry:ex:signxyz}
Find the signature of $xy+yz+zx$ and an orthogonal basis for it.
\end{exercise}

\begin{exercise}[$\star$]\label{geometry:ex:hessian}
Let $f:\mathbb R^n\to\mathbb R$ be of class $C^2$ with a critical point $a$ and Hessian matrix $H$ at $a$. Show that if $H$ is positive definite, $a$ is a
strict local minimum, and if $H$ has signature $(p,m)$ with $p,m\ge1$, $a$ is neither a local minimum nor a local maximum.
\end{exercise}

\begin{exercise}[$\star\star$]\label{geometry:ex:sylvestercrit}
(Sylvester's criterion) Show that a real symmetric matrix is positive definite if and only if all its leading principal minors are positive.
\end{exercise}

\begin{exercise}[$\star$]\label{geometry:ex:binarydisc}
Show that the real form $ax^2+2bxy+cy^2$ is definite if $b^2<ac$, indefinite if $b^2>ac$ and degenerate if $b^2=ac$, and relate this to the roots
of $at^2+2bt+c$.
\end{exercise}

\begin{exercise}[$\star$]\label{geometry:ex:gram}
Let $v_1,\dots,v_k\in\mathbb R^n$ and $G=(v_i\cdot v_j)$ their Gram matrix. Show that $G$ is positive semidefinite, of the same rank as the family $v_i$, and
that $\det G$ is the square of the $k$-dimensional volume of the parallelotope they span.
\end{exercise}

\begin{exercise}[$\star\star$]\label{geometry:ex:lorentz}
On $\mathbb R^4$ let $q(t,x,y,z)=t^2-x^2-y^2-z^2$. Show that if $q(u)>0$ then $\beta(u,v)^2\ge q(u)q(v)$ for every $v$, with equality only if $v$ is a
multiple of $u$: the Cauchy--Schwarz inequality is reversed.
\end{exercise}

\begin{exercise}[$\star$]\label{geometry:ex:conicsclass}
Using Sylvester's law, list all real projective conics, degenerate ones included, up to projective equivalence.
\end{exercise}

\begin{exercise}[$\star$]\label{geometry:ex:saddlelines}
Find all lines on the saddle $z=xy$ through a given point $(a,b,ab)$.
\end{exercise}

\begin{exercise}[$\star$]\label{geometry:ex:wittFp}
Find the Witt index of $x^2+y^2+z^2$ over $\mathbb R$, over $\mathbb Q$ and over $\mathbb F_p$ for odd $p$. Find an isotropic vector for $p=3,5,7$.
\end{exercise}

\begin{exercise}[$\star$]\label{geometry:ex:conicFq}
Show that a nondegenerate conic over $\mathbb F_q$, $q$ odd, has exactly $q+1$ points.
\end{exercise}

\begin{exercise}[$\star\star$]\label{geometry:ex:linesFq}
Show that the hyperbolic quadric over $\mathbb F_q$ contains exactly $2(q+1)$ lines and the elliptic quadric contains none.
\end{exercise}

\begin{exercise}[$\star\star\star$]\label{geometry:ex:skewlines}
Let $L_1,L_2,L_3$ be pairwise skew lines in $\mathbb P^3(K)$, $K$ infinite. Show that exactly one quadric surface contains them, that it is
nondegenerate, and that the lines meeting all three are exactly the lines of its other family.
\end{exercise}

\section{Hints}

\begin{hint}{geometry:ex:signxyz}
Put $x=u+v$, $y=u-v$.
\end{hint}
\begin{hint}{geometry:ex:hessian}
Taylor's formula $f(a+h)=f(a)+\frac12h^{\mathsf T}Hh+o(|h|^2)$, and the smallest eigenvalue of $H$.
\end{hint}
\begin{hint}{geometry:ex:sylvestercrit}
Induction on the size: if the top-left block is positive definite, the signature has $p\ge n-1$. Then look at the sign of the determinant.
\end{hint}
\begin{hint}{geometry:ex:binarydisc}
The determinant of the Gram matrix is $ac-b^2$, the product of the two eigenvalues.
\end{hint}
\begin{hint}{geometry:ex:gram}
$G=M^{\mathsf T}M$ where $M$ has columns $v_i$; for the volume, use an orthonormal basis of the span.
\end{hint}
\begin{hint}{geometry:ex:lorentz}
The orthogonal complement of $u$ is negative definite. Write $v=\lambda u+w$ with $w\perp u$.
\end{hint}
\begin{hint}{geometry:ex:conicsclass}
Conics are forms up to a nonzero factor, so the signature counts only up to the swap $(p,m)\leftrightarrow(m,p)$.
\end{hint}
\begin{hint}{geometry:ex:saddlelines}
A line through the point is $(a+su,\,b+sv,\,ab+sw)$; substitute and require the equation to hold for every $s$.
\end{hint}
\begin{hint}{geometry:ex:wittFp}
A form of dimension $3$ has Witt index at most $1$.
\end{hint}
\begin{hint}{geometry:ex:conicFq}
Proposition~\ref{geometry:prop:chevalley} and the normal form of Appendix~D.
\end{hint}
\begin{hint}{geometry:ex:linesFq}
For the elliptic quadric, look at the tangent plane at a point $P$: it is $P^\perp$. Use Witt cancellation to identify the form on it.
\end{hint}
\begin{hint}{geometry:ex:skewlines}
Quadrics in $\mathbb P^3$ form a space of dimension $9$; a quadric through three points of a line contains the line.
\end{hint}

\section{Solutions}

\begin{solution}{geometry:ex:signxyz}
\idea Turn a product into a difference of squares, as in $xy=\frac14((x+y)^2-(x-y)^2)$.
\writeup With $x=u+v$, $y=u-v$: $xy=u^2-v^2$ and $yz+zx=z(x+y)=2uz$. So the form is $u^2+2uz-v^2=(u+z)^2-z^2-v^2$. Its signature is $(1,2)$. In the
original coordinates, the vectors $e_1=(1,1,0)$, $e_2=(1,-1,0)$, $e_3=(-1,-1,1)$ (corresponding to $(u,v,z)=(1,0,0),(0,1,0),(-1,0,1)$) give
$q(e_1)=1$, $q(e_2)=-1$, $q(e_3)=1-1-1=-1$, and one checks $\beta(e_i,e_j)=0$ for $i\ne j$. As a check with Volume~2: the matrix is
$\frac12(J-I)$, with $J$ the all-ones matrix, whose eigenvalues are $1,-\frac12,-\frac12$.
\end{solution}

\begin{solution}{geometry:ex:hessian}
\idea Near a critical point a function looks like its quadratic form.
\writeup If $H$ is positive definite with smallest eigenvalue $\lambda>0$, then $h^{\mathsf T}Hh\ge\lambda|h|^2$, and
$f(a+h)-f(a)\ge\frac\lambda2|h|^2+o(|h|^2)>0$ for small $h\ne0$. If $p,m\ge1$, choose $h_+$ with $h_+^{\mathsf T}Hh_+>0$ and $h_-$ with
$h_-^{\mathsf T}Hh_-<0$; along $a+th_\pm$ we have $f(a+th_\pm)-f(a)=\frac{t^2}2h_\pm^{\mathsf T}Hh_\pm+o(t^2)$, which takes both signs for small $t$.
The degenerate cases (some zero eigenvalue and no opposite signs) are not decided by $H$: compare $x^2+y^4$ and $x^2-y^4$.
\end{solution}

\begin{solution}{geometry:ex:sylvestercrit}
\idea Positive definiteness passes to subspaces; and a determinant controls the parity of the number of negative eigenvalues.
\writeup If $A$ is positive definite, so is its restriction to the span of the first $k$ basis vectors, whose matrix is the top-left $k\times k$
block $A_k$; its eigenvalues are positive, so $\det A_k>0$. Conversely, suppose all leading minors are positive, and by induction that $A_{n-1}$ is
positive definite. Then $q$ is positive definite on a subspace of dimension $n-1$, so $p\ge n-1$ by Theorem~\ref{geometry:thm:sylvester}. The determinant of $A$
is the product of its eigenvalues and is positive, so it is not zero and the number of negative eigenvalues is even; it is at most $1$, so it is $0$,
and $A$ is positive definite.
\end{solution}

\begin{solution}{geometry:ex:binarydisc}
\idea Two eigenvalues, product $ac-b^2$.
\writeup The Gram matrix $\begin{pmatrix}a&b\\b&c\end{pmatrix}$ has eigenvalues $\lambda_1,\lambda_2$ with $\lambda_1\lambda_2=ac-b^2$. If $b^2<ac$ they
have the same sign (definite); if $b^2>ac$, opposite signs (indefinite, signature $(1,1)$); if $b^2=ac$, one is zero. The form vanishes at $(t,1)$
exactly when $at^2+2bt+c=0$: an indefinite form has two isotropic lines, given by the two real roots, a definite one has none, and a degenerate
nonzero one has a double root.
\end{solution}

\begin{solution}{geometry:ex:gram}
\idea A Gram matrix is $M^{\mathsf T}M$.
\writeup With $M$ the $n\times k$ matrix of columns $v_i$, $G=M^{\mathsf T}M$ and $x^{\mathsf T}Gx=|Mx|^2\ge0$. The kernel of $G$ equals the kernel of $M$
(if $Gx=0$ then $|Mx|^2=x^{\mathsf T}Gx=0$), so they have the same rank. In an orthonormal basis of a $k$-dimensional space containing the $v_i$, $M$
becomes a $k\times k$ matrix $M'$ with the same Gram matrix, and $\det G=(\det M')^2$, the square of the volume (Volume~2, Chapter~2).
\end{solution}

\begin{solution}{geometry:ex:lorentz}
\idea In relativity this is the twin paradox: the straight path between two events is the one with the \emph{longest} proper time.
\writeup The signature of $q$ is $(1,3)$. Since $q(u)>0$, $V=\mathbb Ru\perp u^\perp$, and the signature of $u^\perp$ is $(0,3)$ by Sylvester's law: $q$ is
negative definite on $u^\perp$. Write $v=\lambda u+w$ with $w\in u^\perp$. Then $\beta(u,v)=\lambda q(u)$ and $q(v)=\lambda^2q(u)+q(w)\le\lambda^2q(u)$. Hence
$\beta(u,v)^2=\lambda^2q(u)^2\ge q(u)q(v)$, with equality if and only if $q(w)=0$, that is $w=0$.
\end{solution}

\begin{solution}{geometry:ex:conicsclass}
\idea Classify $3\times3$ real symmetric matrices up to congruence and a nonzero factor.
\writeup Up to swapping $p$ and $m$, the signatures of nonzero forms in three variables are: $(3,0)$, the empty conic $x^2+y^2+z^2=0$; $(2,1)$, the circle
$x^2+y^2=z^2$; $(2,0)$, the single point $x^2+y^2=0$, which is $[0:0:1]$; $(1,1)$, two lines $x^2-y^2=0$; $(1,0)$, the double line $x^2=0$. So there are
five types, two of them nondegenerate.
\end{solution}

\begin{solution}{geometry:ex:saddlelines}
\idea A line lies on the surface when a polynomial in $s$ of degree two vanishes identically.
\writeup Substituting $(a+su,b+sv,ab+sw)$ into $z=xy$ gives $ab+sw=ab+s(av+bu)+s^2uv$, so we need $uv=0$ and $w=av+bu$. If $u=0$, the direction is
$(0,1,a)$: the line $x=a$, $z=ay$. If $v=0$, it is $(1,0,b)$: the line $y=b$, $z=bx$. So exactly two lines pass through each point, one from each family.
\end{solution}

\begin{solution}{geometry:ex:wittFp}
\idea Isotropy is all that matters in dimension three.
\writeup A $3$-dimensional form cannot contain $2\mathbb H$, so its Witt index is $0$ or $1$ according as it is anisotropic or not. Over $\mathbb R$ and
over $\mathbb Q$, $x^2+y^2+z^2$ is positive definite, hence anisotropic: index $0$. Over $\mathbb F_p$ it is isotropic by Proposition~\ref{geometry:prop:chevalley}:
index $1$. Isotropic vectors: $(1,1,1)$ for $p=3$; $(1,2,0)$ for $p=5$, since $1+4=5$; $(1,2,3)$ for $p=7$, since $1+4+9=14$.
\end{solution}

\begin{solution}{geometry:ex:conicFq}
\idea An isotropic vector is a point, and a conic with a point is a projective line.
\writeup The ternary form of the conic is nondegenerate, so it is isotropic by Proposition~\ref{geometry:prop:chevalley}: the conic has a point. By the normal
form of Theorem~\ref{geometry:thm:normalform}, it is projectively equivalent to $y^2=xz$, the image of the bijection $\varphi$ from $\mathbb P^1(\mathbb F_q)$,
which has $q+1$ points.
\end{solution}

\begin{solution}{geometry:ex:linesFq}
\idea Lines on the hyperbolic quadric are the two rulings; on the elliptic quadric the tangent plane at a point meets it only at that point.
\writeup The hyperbolic quadric is $x_0x_3=x_1x_2$, and by Theorem~\ref{geometry:thm:segre} its lines are the $q+1$ lines of each family, $2(q+1)$ in all. For the
elliptic quadric $\mathbb H\perp N$, let $P=[p]$ be a point. A line through $P$ on the quadric lies in the tangent plane $p^\perp$ (if $[x]$ and $[p]$ and
the whole line are on the quadric, then $0=q(p+x)=2\beta(p,x)$). By Lemma~\ref{geometry:lem:hyperbolic}, $p$ lies in a hyperbolic plane $H'$, and
$V=H'\perp H'^\perp$. By Witt cancellation, $\mathbb H\perp N\cong H'\perp H'^\perp$ gives $H'^\perp\cong N$, anisotropic. Now $p^\perp=Kp\oplus H'^\perp$, and for
$x=\lambda p+w$ in it, $q(x)=q(w)$, which vanishes only for $w=0$. So the only point of the quadric in the tangent plane is $P$, and no line of the
quadric passes through $P$.
\end{solution}

\begin{solution}{geometry:ex:skewlines}
\idea Nine points impose at most nine conditions on a ten-dimensional space of quadratic forms; skew lines force the quadric to be the Segre quadric,
and the other ruling forces uniqueness.
\writeup \emph{Existence.} Quadratic forms in four variables form a space of dimension $10$. Choose three points on each line; requiring a form to vanish
at these nine points is nine linear conditions, so a nonzero form $f$ vanishes at all of them. On each line $L_i$ the restriction of $f$ is a binary
quadratic form with three roots, so it is zero: the quadric $\{f=0\}$ contains the three lines.

\emph{Nondegeneracy.} Over $\overline K$, a degenerate quadric is a cone with a vertex (rank $3$), a pair of planes (rank $2$) or a double plane (rank $1$).
Every line on a cone passes through its vertex, so two of our lines would meet. Three lines in the union of two planes put two of them in the same plane,
where they meet. So $f$ is nondegenerate. It contains a line, that is, a two-dimensional subspace $U$ on which the form vanishes. Take $e\in U$ nonzero and a
hyperbolic plane $H=\langle e,f'\rangle$ as in Lemma~\ref{geometry:lem:hyperbolic}. The space $U\cap f'^\perp$ is at least one-dimensional; a nonzero $u$ in it is
orthogonal to $f'$ and (being in $U$) to $e$, so $u\in H^\perp$ is isotropic. So $H^\perp$ is a hyperbolic plane too, the Witt index is $2$, and after a
change of coordinates $f=x_0x_3-x_1x_2$. The three lines are pairwise disjoint, so they all belong to one family of Theorem~\ref{geometry:thm:segre}.

\emph{The other family and uniqueness.} A line $M$ of the other family meets each $L_i$ in one point, and these three points are distinct (the $L_i$ are
disjoint). Conversely, a line meeting $L_1,L_2,L_3$ meets the quadric in three points, so it lies on it, and since it meets $L_1$ it is not in the first family.
Now let $g$ be another quadratic form vanishing on the three lines. Each line $M$ of the other family contains three zeros of $g$, so $g$ vanishes on $M$, hence
on all of $Q$. Then $g(su,sv,tu,tv)=0$ as a polynomial in $s,t,u,v$ ($K$ is infinite). Of the ten monomials of degree two, only $x_0x_3$ and $x_1x_2$ give the
same monomial $stuv$, so all the other coefficients of $g$ vanish and the coefficients of $x_0x_3$ and $x_1x_2$ are opposite: $g$ is a multiple of $f$.
\end{solution}

\section*{Sources and further reading}

\begin{itemize}
\item J.~J.~Sylvester, \emph{A demonstration of the theorem that every homogeneous quadratic polynomial is reducible by real orthogonal
substitutions\dots}, Philosophical Magazine (4) 4 (1852), 138--142 (the epigraph is the beginning of its long title).
\item R.~Graefe et al., \emph{Vladimir G.~Suchov 1853--1939: Die Kunst der sparsamen Konstruktion}, Deutsche
Verlags-Anstalt, 1990 (Shukhov's hyperboloid towers, the 1896 Nizhny Novgorod tower and the Shabolovka tower); \emph{The Nijni-Novgorod exhibition:
water tower}, The Engineer, 19 March 1897, 292--294.
\item E.~Witt, \emph{Theorie der quadratischen Formen in beliebigen K\"orpern}, Journal f\"ur die reine und angewandte Mathematik 176 (1937), 31--44
(cancellation and the Witt decomposition).
\item T.~Y.~Lam, \emph{Introduction to Quadratic Forms over Fields}, AMS, 2005; J.-P.~Serre, \emph{A Course in Arithmetic}, Springer, 1973, Chapter~IV
(Hasse--Minkowski); J.~W.~S.~Cassels, \emph{Rational Quadratic Forms}, Academic Press, 1978.
\item M.~Berger, \emph{Geometry II}, Springer, 1987, Chapters~13--15 (projective and affine quadrics, the seventeen real types); D.~Hilbert and
S.~Cohn-Vossen, \emph{Geometry and the Imagination}, Chelsea, 1952, Chapter~I (ruled surfaces and Shukhov-style models).
\item J.~W.~P.~Hirschfeld, \emph{Finite Projective Spaces of Three Dimensions}, Oxford, 1985 (quadrics over finite fields, ovoids).
\end{itemize}
